\scanpage{55}
\mspage{11}

3) \deleted{Gommage} \inserted{Bouchage de trous, ou ``gommage''} de cercles de recollement.

Cette opération est relative à la donnée d'une $1$-sous-variété fermée de $\partial X$, i.e. d'une partie $A'_l$ de $\pi_0(\partial X)\simeq A_l$ (ens des arêtes libres du graphe). Elle consiste à boucher les trous correspondants, au moyen \unsure{des|de} cônes sur les composantes connexes du bord de $\Sigma'_l$. Au niveau du graphe, cela consiste à gommer les arêtes libres correspondantes, avec bien entendu aussi les repères de marquage qui se trouvent au dessus --- \deleted{mais} les pondérations des \unsure{sommet|sommets} restent inchangées.

\deleted{\unsure{4)}} \uline{Gommage simultané mixte}, \deleted{\unsure{au}} correspondant à une partie \deleted{$A'\amalg R'$} \inserted{$R'\amalg\Sigma'$} \deleted{quelconque} de $R\amalg\Sigma$ \deleted{\unsure{du}} \deleted{\unsure{de} $A\amalg R$,} \inserted{qui soit à la fois ouverte et fermée} \deleted{i.e.} \inserted{et telle que $R'\supset R\cap\Sigma'$,} \deleted{une \unsure{ens.} de composantes connexes de $R\amalg\Sigma$,} ou ce qui revient au $=$, une partie \inserted{$R'\amalg A'$} de $R\amalg A$, telle que $R'\supset R|A'$. Au niveau du graphe pondéré à repères, cette opération consiste à enlever \deleted{\unsure{les}} \inserted{de $\vec R$} la partie corr. $\vec R'$, à contracter en un point les arêtes de $\unsure{A'|\Pi'}$\deleted{\unsure{pas}} (aussi bien libres que liées), en faisant suivre \unsure{les pondérations|la pondération}.
% archive sheet number: 54

